% 图契约：固定回路的感生电动势与运动导体的动生电动势；几何、方向和公式与正文同一约定。
\documentclass[tikz,border=5pt]{standalone}
\usepackage{fontspec,xeCJK,amsmath,amssymb}
\IfFontExistsTF{Times New Roman}{\setmainfont{Times New Roman}}{\setmainfont{TeX Gyre Termes}}
\IfFontExistsTF{Heiti SC}{\setCJKmainfont[AutoFakeBold=1.5]{Heiti SC}}{\setCJKmainfont[AutoFakeBold=1.5]{FandolHei-Regular}}
\usetikzlibrary{arrows.meta,calc,patterns}
\definecolor{ink}{HTML}{243745}
\definecolor{blue}{HTML}{275DAB}
\definecolor{green}{HTML}{177C63}
\definecolor{red}{HTML}{BD4935}
\definecolor{orange}{HTML}{AA6C16}
\definecolor{purple}{HTML}{715393}
\begin{document}
\begin{tikzpicture}[x=1cm,y=1cm,font=\fontsize{9.5}{12}\selectfont,
v/.style={-{Latex[length=2mm,width=1.3mm]},line width=.85pt},
axis/.style={-{Latex[length=1.7mm,width=1.1mm]},ink!65,line width=.55pt},
aux/.style={ink!45,dashed,line width=.55pt},
lab/.style={fill=white,inner sep=1.8pt},
card/.style={draw=ink!25,fill=blue!3,rounded corners=2pt,text width=6.1cm,minimum height=.85cm,inner sep=6pt,align=center}]
\path[use as bounding box](-8,-3.6)rectangle(8,3.65);
\node[font=\bfseries\fontsize{10.5}{13}\selectfont,text width=15.1cm,align=center]at(0.00000,3.22000){M7-H17　固定回路的感生电动势与运动导体的动生电动势};
\begin{scope}[shift={(-4,0)}]
\draw[green](-2,-1.3)circle(0.13);
\fill[green](-2,-1.3)circle(.035);
\draw[green](-2,0)circle(0.13);
\fill[green](-2,0)circle(.035);
\draw[green](-2,1.3)circle(0.13);
\fill[green](-2,1.3)circle(.035);
\draw[green](-1,-1.3)circle(0.13);
\fill[green](-1,-1.3)circle(.035);
\draw[green](-1,0)circle(0.13);
\fill[green](-1,0)circle(.035);
\draw[green](-1,1.3)circle(0.13);
\fill[green](-1,1.3)circle(.035);
\draw[green](0,-1.3)circle(0.13);
\fill[green](0,-1.3)circle(.035);
\draw[green](0,0)circle(0.13);
\fill[green](0,0)circle(.035);
\draw[green](0,1.3)circle(0.13);
\fill[green](0,1.3)circle(.035);
\draw[green](1,-1.3)circle(0.13);
\fill[green](1,-1.3)circle(.035);
\draw[green](1,0)circle(0.13);
\fill[green](1,0)circle(.035);
\draw[green](1,1.3)circle(0.13);
\fill[green](1,1.3)circle(.035);
\draw[green](2,-1.3)circle(0.13);
\fill[green](2,-1.3)circle(.035);
\draw[green](2,0)circle(0.13);
\fill[green](2,0)circle(.035);
\draw[green](2,1.3)circle(0.13);
\fill[green](2,1.3)circle(.035);
\draw[ink,line width=2pt](0.00000,-1.80000)--(0.00000,1.80000);
\node[lab]at(-0.30000,-1.90000){$A,+$};
\node[lab]at(-0.30000,1.90000){$B,-$};
\draw[v,blue](0.40000,0.80000)--(2.00000,0.80000);
\node[lab,text=blue]at(1.30000,1.15000){$\mathbf v$};
\draw[v,red](-0.45000,0.70000)--(-0.45000,-0.60000);
\node[lab,text=red]at(-1.35000,0.30000){$\mathbf v\times\mathbf B$};
\draw[{Latex[length=1.3mm]}-{Latex[length=1.3mm]},ink!65,line width=.5pt](2.70000,-1.80000)--(2.70000,1.80000);
\node[lab]at(3.00000,0.00000){$\ell$};
\end{scope}
\node[card]at(4.00000,1.25000){$\mathcal E=\oint(\mathbf E+\mathbf v\times\mathbf B)\cdot d\mathbf l$};
\node[card]at(4.00000,-1.05000){$\mathcal E=-d\Phi_B/dt$};
\end{tikzpicture}\end{document}
